% critique.tex — f171v, vv. 764–800 : commentaire critique
% Articles textuels tirés du commentaire de l'auteur (latex/Aristophane-Thesmo Harmattan 2023 vv 642-fin-bis.tex),
% texte et notes conservés ; répartition détaillée dans repartition.md.
\cc{772}{πόθεν οὖν γένοιντ᾽ ἄν μοι πλάται. πόθεν;}{Le vers est incomplet. Scaliger proposa la solution la plus instinctive : redoubler le πόθεν en fin de vers : πόθεν ; <πόθεν;> Cette solution est vraisemblablement la bonne. \ \ \ 1° \ Une simple haplographie l'explique. \ \ \ 2° \ Elle est une nouvelle occurrence des railleries aristophanesques quant à la manie du redoublement chez les tragiques. Le double πόθεν se trouve d'ailleurs dans l'\italique{Ion} (vers 1453 sq.) :\citComm{ἰὼ γύναι, πόθεν πόθεν ἔλαβες ἐμὸν\\
βρέφος ἐς ἀγκάλας;}Même si ce qui faisait vraiment battre le cœur d'Euripide, c'était le redoublement du ποῦ (\italique{Électre}, v. 487 sq.) :\citComm{ποῦ ποῦ νεᾶνις πότνι᾽ ἐμὴ δέσποινά τε,\\
Ἀγαμέμνονος παῖς, ὅν ποτ᾽ ἐξέθρεψ᾽ ἐγώ;}Dans les \italique{Oiseaux} (vers 921), c'est le redoublement de πάλαι qui est moqué : l'ennuyeux poète qui débite\citComm{πάλαι, πάλαι δὴ τήνδ᾽ ἐγὼ κλῄζω πόλιν}ne pénétrera pas plus dans Νεφελοκοκκυγίᾳ qu'il n'aura de chance de demeurer dans la Cité de Platon.}
\cc{781}{τουτὶ τὸ ῥῶ μοχθηρόν·}{« Voilà un `Rhô' défectueux ! »... telle est la leçon de notre unique manuscrit et telles aussi les traductions embarrassées. Car la question fuse : pourquoi Mnésiloque nous parle-t-il d'un `ρ' ? Le scholiaste toujours prompt à lancer une ânerie propose d'y voir la troisième lettre du destinataire, Εὐρῑπίδη... ce qui ne répond pas à la question : pourquoi la troisième lettre ? Passons sur les autres explications du même acabit (si Aristophane avait un peu plus de glosateurs, tous les mots du dictionnaire comportant un rhô y seraient effectivement passés : Εὐριπίδη, χώρει, χαίρειν, \&c.)
Lisez simplement : \citComm{τουτὶ τορῶ μοχθηρόν.}L'aoriste τορεῖν est homérique et fait partie d'un groupe dialectal, également homérique, d'aoristes radicaux thématiques où s'est dégagé un vocalisme en `o' : ἔμολον de βλώσκω, ἔθορον de θρώσκω et ἔτορον de τείρω. Ces homérismes se sont introduits chez les Tragiques, et peuvent aussi se rencontrer dans certains passages d'Aristophane. On l'avait déjà noté pour ἔμολον : on peut maintenant y ajouter ἔτορον\nc{Je dois cette remarque à Jean-Victor Vernhes qui, entre autres qualités, est un très fin connaisseur de la langue homérique. Que le lecteur n'hésite pas à consulter le Dictionnaire étymologique de la langue grecque de Pierre Chantraine à l'entrée τορεῖν pour plus de détails.}. Hésychius (τ 1162) -- qui me fut d'un grand secours comme à son habitude – l'avait déjà relevé :\citComm{τορεῖν· τορῆσαι, τρῆσαι, τεμεῖν ;}le verbe signifie \italique{percer}, \italique{perforer}, \italique{couper}.}
\cc{782}{ποίαν αὔλακα}{Pour le sens de ποίαν (reprise d'un terme), voir le commentaire de la pièce. Et si l'on en doutait, l'on pourrait toujours corriger le texte -- ce dont on se gardera bien ! Lire ποῖ ποτ' ἂν αὔλακα revient au même quant à l'action et dispense de supposer que l'αὖλαξ était le terme employé par Euripide. (ποτε \italique{gratia metri} mais dont la disparition relèverait d'une haplographie bien explicable. Au reste, la collocation `ποῖ ποτε' est très fréquente.)}
\cc{797}{τὸ κακὸν ζητεῖτε θεᾶσθαι}{Suivant la belle correction de Kaye : τὸ κακὸν ζητεῖτε θεᾶσθαι, v. 797, au lieu de : ζητεῖ τὸ κακὸν τε θεᾶσθαι. (Note de Willems.)}
